\chapter{Spinor Fields, Fermionic Statistics, and Grassmann Path Integrals}
\label{ch:spinor-fields-fermionic-path-integrals}

\section{The Free Dirac Field}

Work in \(D=4\) with
\[
  \eta_{\mu\nu}=\operatorname{diag}(-,+,+,+),
  \qquad
  \{\gamma^\mu,\gamma^\nu\}=2\eta^{\mu\nu}.
\]
Set
\[
  \not\partial=\gamma^\mu\partial_\mu,
  \qquad
  \not p=p_\mu\gamma^\mu,
  \qquad
  \beta=\ii\gamma^0 .
\]
For a spinor \(s\), its Dirac adjoint is
\[
  \bar s=s^\dagger\beta .
\]

Let \(u^\sigma(p)\) and \(v^\sigma(p)\), \(\sigma=1,2\), be the spinor
intertwiners on the positive-energy mass shell \(\Sigma_m^+\) constructed in
the previous chapter:
\[
  (\not p-\ii m)u^\sigma(p)=0,
  \qquad
  (\not p+\ii m)v^\sigma(p)=0.
\]
Let
\[
  \dd\mu_m(p)
  =
  \frac{\dd^3\vec p}{(2\pi)^3\,2\omega_{\vec p}},
  \qquad
  \omega_{\vec p}=\sqrt{\vec p^{\,2}+m^2},
  \qquad
  p^0=\omega_{\vec p}.
\]
Introduce two independent families of fermionic creation and annihilation
operators
\[
  b_\sigma(\vec p),\ b_\sigma^\dagger(\vec p),
  \qquad
  d_\sigma(\vec p),\ d_\sigma^\dagger(\vec p),
\]
with canonical anticommutation relations
\[
  \{b_\sigma(\vec p),b_\tau^\dagger(\vec q)\}
  =
  (2\pi)^3\,2\omega_{\vec p}\,
  \delta_{\sigma\tau}\delta^3(\vec p-\vec q),
\]
\[
  \{d_\sigma(\vec p),d_\tau^\dagger(\vec q)\}
  =
  (2\pi)^3\,2\omega_{\vec p}\,
  \delta_{\sigma\tau}\delta^3(\vec p-\vec q),
\]
all other anticommutators being zero. The vacuum vector \(\Omega\) is
annihilated by all \(b_\sigma(\vec p)\) and \(d_\sigma(\vec p)\).

The charged Dirac field is the operator-valued distribution
\[
  \widehat\psi_\alpha(x)
  =
  \sum_{\sigma=1}^2
  \int_{\Sigma_m^+}\dd\mu_m(p)\,
  \left[
    v_\alpha^\sigma(p)b_\sigma^\dagger(\vec p)\ee^{-\ii p\cdot x}
    +
    u_\alpha^\sigma(p)d_\sigma(\vec p)\ee^{\ii p\cdot x}
  \right].
\]
Its adjoint field is
\[
  \widehat{\bar\psi}^{\alpha}(x)
  =
  \bigl(\widehat\psi^\dagger(x)\beta\bigr)^\alpha .
\]
Because of the two momentum-space Dirac equations above,
\[
  (\not\partial+m)\widehat\psi(x)=0,
  \qquad
  \widehat{\bar\psi}(x)(-\overleftarrow{\not\partial}+m)=0,
\]
where the arrow indicates that the derivative acts on the field to its left.

The convention in this chapter assigns charge \(+1\) to the one-particle
states created by \(b^\dagger\) and charge \(-1\) to the states created by
\(d^\dagger\). If
\[
  U_\theta=\ee^{\ii\theta \widehat Q},
  \qquad
  U_\theta\Omega=\Omega ,
\]
then
\[
  U_\theta\widehat\psi_\alpha(x)U_\theta^{-1}
  =
  \ee^{\ii\theta}\widehat\psi_\alpha(x),
  \qquad
  U_\theta\widehat{\bar\psi}^{\alpha}(x)U_\theta^{-1}
  =
  \ee^{-\ii\theta}\widehat{\bar\psi}^{\alpha}(x).
\]

\section{Locality and the CAR}

Define the positive-frequency scalar distribution
\[
  \Delta_+(x;m^2)
  =
  \int_{\Sigma_m^+}\dd\mu_m(p)\,\ee^{\ii p\cdot x}.
\]
With the spinor normalizations fixed by the one-particle inner product, the
field anticommutator is
\[
  \{\widehat\psi_\alpha(x),\widehat{\bar\psi}^{\beta}(y)\}
  =
  (-\not\partial_x+m)_\alpha{}^\beta
  \left[
    \Delta_+(x-y;m^2)-\Delta_+(y-x;m^2)
  \right].
\]
The bracketed distribution is the Pauli--Jordan distribution up to the
normalization convention already used for the scalar field. It vanishes at
spacelike separation. Hence
\[
  \{\widehat\psi_\alpha(x),\widehat{\bar\psi}^{\beta}(y)\}=0
  \qquad\text{when }(x-y)^2>0.
\]
The equal-time anticommutator is obtained by differentiating the
Pauli--Jordan distribution:
\[
  \{\widehat\psi_\alpha(t,\vec x),
    \widehat{\bar\psi}^{\beta}(t,\vec y)\}
  =
  \ii(\gamma^0)_\alpha{}^\beta\delta^3(\vec x-\vec y),
\]
and
\[
  \{\widehat\psi_\alpha(t,\vec x),\widehat\psi_\beta(t,\vec y)\}=0,
  \qquad
  \{\widehat{\bar\psi}^{\alpha}(t,\vec x),
    \widehat{\bar\psi}^{\beta}(t,\vec y)\}=0.
\]

The local algebra is therefore a graded local algebra. If \(|A|\in\{0,1\}\)
is the fermion parity of a homogeneous operator \(A\), the graded commutator is
\[
  [A,B]_{\mathrm g}
  =
  AB-(-1)^{|A||B|}BA .
\]
Locality is the vanishing of this graded commutator for spacelike separated
smeared fields. The Dirac field has parity \(1\), while bosonic local fields
have parity \(0\).

\begin{remark}
The preceding computation is the free spin-\(\frac12\) instance of the
spin-statistics connection. In a Wightman setting with positivity of energy,
locality, covariance, and the usual domain assumptions, the spin-statistics
theorem identifies integer-spin fields with even local commutation and
half-integer-spin fields with odd local anticommutation. This chapter uses the
theorem only as orientation; the displayed anticommutator is the direct
verification needed for the free Dirac field used below.
\end{remark}

\section{Dirac, Weyl, and Majorana Fields}

The Dirac field transforms in the spinor representation:
\[
  U(\Lambda)\widehat\psi_\alpha(x)U(\Lambda)^{-1}
  =
  R_\alpha{}^\beta(\Lambda)\widehat\psi_\beta(\Lambda x),
\]
where
\[
  R(\Lambda)
  =
  \exp\!\left(-\frac{\ii}{2}\omega_{\mu\nu}S^{\mu\nu}\right),
  \qquad
  S^{\mu\nu}=-\frac{\ii}{4}[\gamma^\mu,\gamma^\nu],
\]
for \(\Lambda=\exp\omega\). The field \(\widehat\psi\) has four complex
components before additional Lorentz-covariant conditions are imposed.

Define
\[
  \gamma_5=-\ii\gamma^0\gamma^1\gamma^2\gamma^3 .
\]
It obeys
\[
  \gamma_5^2=1,
  \qquad
  \{\gamma_5,\gamma^\mu\}=0,
  \qquad
  [\gamma_5,S^{\mu\nu}]=0.
\]
Therefore
\[
  P_\pm=\frac{1\pm\gamma_5}{2}
\]
are Lorentz-covariant projectors. A Weyl spinor field is a Dirac spinor field
subject to
\[
  \widehat\chi_\pm=P_\pm\widehat\chi_\pm .
\]
The two chiralities are inequivalent complex representations of
\(\operatorname{Spin}^+(1,3)\). The Dirac mass term couples them because
\(\gamma^\mu\) maps one chirality to the other.

A Majorana condition is a Lorentz-covariant reality condition. Suppose a
matrix \(B\) is chosen so that
\[
  (\gamma^\mu)^*=B\gamma^\mu B^{-1},
  \qquad
  B^{-1}(B^{-1})^*=1.
\]
Here \( * \) denotes the antilinear operation on spinor components, combined
with componentwise adjunction when it is applied to operator-valued spinor
fields. Then
\[
  \widehat\psi^c=B^{-1}\widehat\psi^*
\]
transforms in the same Lorentz representation as \(\widehat\psi\). A Majorana
field satisfies
\[
  \widehat\xi=\widehat\xi^c .
\]
If \(\widehat\chi_+=P_+\widehat\chi_+\) is a Weyl field, then
\[
  \widehat\xi
  =
  \widehat\chi_+
  +
  B^*\widehat\chi_+^*
\]
is Majorana, and
\[
  P_+\widehat\xi=\widehat\chi_+ .
\]
Thus in four dimensions a Majorana spinor is equivalent, as Lorentz
representation data, to one Weyl spinor together with its conjugate.

\section{Berezin Calculus}

Let \(\Lambda_N\) be the complex Grassmann algebra generated by
\[
  \eta_1,\ldots,\eta_N
\]
with relations
\[
  \eta_a\eta_b=-\eta_b\eta_a,
  \qquad
  \eta_a^2=0.
\]
A homogeneous monomial has parity equal to its degree modulo \(2\). If \(X\)
and \(Y\) are homogeneous, then
\[
  XY=(-1)^{|X||Y|}YX
\]
when every generator appearing in \(X\) is distinct from every generator
appearing in \(Y\).

The left derivative \(\partial^L/\partial\eta_a\) is defined by
\[
  \frac{\partial^L}{\partial\eta_a}\eta_b=\delta_{ab},
  \qquad
  \frac{\partial^L}{\partial\eta_a}(XY)
  =
  \frac{\partial^L X}{\partial\eta_a}Y
  +
  (-1)^{|X|}X\frac{\partial^L Y}{\partial\eta_a}
\]
for homogeneous \(X\). Berezin integration is the same algebraic operation:
\[
  \int\dd\eta_a\,1=0,
  \qquad
  \int\dd\eta_a\,\eta_a=1.
\]
For \(N\) generators the ordered integral is fixed by
\[
  \int\dd\eta_N\cdots\dd\eta_1\,
  \eta_1\eta_2\cdots\eta_N
  =
  1.
\]

For independent pairs \(\eta_a,\bar\eta_a\), \(a=1,\ldots,N\), and an
invertible \(N\times N\) complex matrix \(A\), use the ordered measure
\[
  [\dd\bar\eta\,\dd\eta]
  =
  \prod_{a=1}^N\dd\bar\eta_a\,\dd\eta_a .
\]
With this convention,
\[
  \int[\dd\bar\eta\,\dd\eta]\,
  \exp\left(-\bar\eta_a A^a{}_b\eta_b\right)
  =
  \det A .
\]
Adding Grassmann sources \(J,\bar J\),
\[
\begin{aligned}
  &\int[\dd\bar\eta\,\dd\eta]\,
  \exp\left(
    -\bar\eta_a A^a{}_b\eta_b
    +\bar J\eta+\bar\eta J
  \right)  \\
  &\hspace{2.2cm}
  =
  \det A\,
  \exp\left(\bar J A^{-1}J\right).
\end{aligned}
\]
The normalized two-point contraction is
\[
  \langle\eta_a\bar\eta_b\rangle_A
  =
  (A^{-1})_a{}^b .
\]
Higher Gaussian moments are obtained by repeated contractions, with the sign
determined by the number of odd variables moved past one another to bring a
contracted pair together.

\section{The Fermionic Coherent-State Path Integral}

The elementary complex fermionic system has two operators
\[
  \widehat\eta,\qquad \widehat\chi,
\]
obeying
\[
  \{\widehat\eta,\widehat\chi\}=1,
  \qquad
  \widehat\eta^2=\widehat\chi^2=0.
\]
On the two-dimensional Hilbert space with basis \(\ket{0},\ket{1}\), choose
\[
  \widehat\eta\ket{0}=0,\qquad
  \widehat\eta\ket{1}=\ket{0},
  \qquad
  \widehat\chi\ket{0}=\ket{1},\qquad
  \widehat\chi\ket{1}=0.
\]

\begin{center}
\begin{tikzpicture}[>=Stealth,scale=0.95]
  \node[gray!70] (nullL) at (-1.25,0) {\(0\)};
  \node[draw,circle,minimum size=0.95cm] (zero) at (0,0) {\(\ket{0}\)};
  \node[draw,circle,minimum size=0.95cm] (one) at (4,0) {\(\ket{1}\)};
  \node[gray!70] (nullR) at (5.25,0) {\(0\)};
  \draw[->,thick,blue!70!black]
    (zero) to[bend left=28] node[above] {\(\widehat\chi\)} (one);
  \draw[->,thick,green!50!black]
    (one) to[bend left=28] node[below] {\(\widehat\eta\)} (zero);
  \draw[->,gray!65] (zero) -- (nullL)
    node[midway,above] {\(\widehat\eta\)};
  \draw[->,gray!65] (one) -- (nullR)
    node[midway,above] {\(\widehat\chi\)};
\end{tikzpicture}
\end{center}

For a Grassmann parameter \(\eta\), define the right coherent ket
\[
  \ket{\eta}=\ket{0}+\ket{1}\eta ,
\]
so that
\[
  \widehat\eta\ket{\eta}=\ket{\eta}\eta .
\]
Define the dual coherent bra by
\[
  \langle\!\langle\chi|
  =
  -\bra{1}+\chi\bra{0}.
\]
It satisfies
\[
  \langle\!\langle\chi|\eta\rangle=\chi-\eta .
\]
The resolution of the identity is
\[
  \int\ket{\eta}\,\dd\eta\,\langle\!\langle\eta|
  =
  \mathbf 1 .
\]

Let
\[
  U(T)=\ee^{-\ii \widehat H T}.
\]
Insert the coherent-state resolution of identity at times
\[
  0=t_0<t_1<\cdots<t_N=T,
  \qquad
  t_{n+1}-t_n=\varepsilon=T/N.
\]
With \(\eta_0=\eta_i\) and \(\eta_N=\eta_f\), the time-sliced kernel is
\[
\begin{aligned}
  \langle\!\langle\eta_f|U(T)|\eta_i\rangle
  =
  \lim_{N\to\infty}
  \int
  \prod_{n=0}^{N-1}\dd\chi_n
  \prod_{n=1}^{N-1}\dd\eta_n\,
  \exp
  \sum_{n=0}^{N-1}
  \left[
    -\chi_n(\eta_{n+1}-\eta_n)
    -\ii\varepsilon H(\chi_n,\eta_n)
  \right].
\end{aligned}
\]
The continuum notation for the same construction is
\[
  \int D\chi\,D\eta\,
  \exp\left\{
    \int_0^T\dd t\,
    \left[-\chi\dot\eta-\ii H(\chi,\eta)\right]
  \right\}.
\]
Equivalently, with the first-order action
\[
  S[\chi,\eta]
  =
  \int_0^T\dd t\,\left[\ii\chi\dot\eta-H(\chi,\eta)\right],
\]
the continuum Grassmann path-integral notation assigns the weight
\(\ee^{\ii S}\) to the kernel. This expression denotes the limit, when it exists,
of the finite time-sliced Berezin integrals displayed above; the time-sliced
expression is the regulated meaning of the notation.

\section{Dirac Functional Integrals and Propagators}

For a classical Grassmann spinor field \(\psi_\alpha(x)\), define
\[
  \bar\psi^\alpha(x)
  =
  \bigl(\psi^\dagger(x)\beta\bigr)^\alpha .
\]
The free Dirac Lagrangian density in the conventions of this volume is
\[
  \mathcal L_{\mathrm D}
  =
  -\bar\psi(\not\partial+m)\psi .
\]
Its Euler--Lagrange equation is
\[
  (\not\partial+m)\psi=0.
\]
The formal Lorentzian generating functional is normalized so that
\[
\begin{aligned}
  Z[\bar J,J]
  &=
  \int D\bar\psi\,D\psi\,
  \exp\left[
    \ii\int\dd^4x\,\mathcal L_{\mathrm D}
    +
    \ii\int\dd^4x\,
    \bigl(\bar J\psi+\bar\psi J\bigr)
  \right]  \\
  &=
  Z[0,0]\,
  \exp\left[
    -\ii\int\dd^4x\,\dd^4y\,
    \bar J^\alpha(x)
    S_{F,\alpha}{}^\beta(x-y)
    J_\beta(y)
  \right],
\end{aligned}
\]
where \(J\) and \(\bar J\) are Grassmann sources. The Feynman two-point
function is
\[
  S_{F,\alpha}{}^\beta(x-y)
  =
  \langle\Omega|
  \operatorname T\{
    \widehat\psi_\alpha(x)\widehat{\bar\psi}^{\beta}(y)
  \}
  |\Omega\rangle .
\]
It is the Feynman inverse of \(\not\partial+m\):
\[
  S_F(x;m)
  =
  (-\not\partial_x+m)\Delta_F(x;m^2).
\]
Using the scalar convention
\[
  \Delta_F(x;m^2)
  =
  -\ii
  \int\frac{\dd^4k}{(2\pi)^4}
  \frac{\ee^{\ii k\cdot x}}{k^2+m^2-\ii\epsilon},
\]
one obtains
\[
  S_F(k)
  =
  \frac{-\ii(-\ii\not k+m)}
       {k^2+m^2-\ii\epsilon}.
\]
The numerator is the spinor projector algebra of the previous chapter, while
the denominator is the scalar mass-shell pole with the Feynman boundary value.

\section{Spinor Spectral Poles and LSZ}

Let an interacting theory have a global \(U(1)\) charge \(\widehat Q\), a
vacuum satisfying \(\widehat Q\Omega=0\), and an isolated massive
spin-\(\frac12\) one-particle sector of mass \(m\) and charge \(+1\). Let
\(\widehat\psi_\alpha\) be a local spinor field of charge \(+1\) with nonzero
overlap with this sector. After choosing the phase of one-particle states,
write
\[
  \langle\vec p,\sigma,+|\widehat\psi_\alpha(x)|\Omega\rangle
  =
  Z_\psi^{1/2}\,
  \ee^{-\ii p\cdot x}v_\alpha^\sigma(p),
  \qquad
  Z_\psi\ge0 .
\]
The conjugate field has the corresponding overlap with the charge \(-1\)
sector,
\[
  \langle\vec p,\sigma,-|
  \widehat{\bar\psi}^{\alpha}(x)|\Omega\rangle
  =
  Z_\psi^{1/2}\,
  \ee^{-\ii p\cdot x}\bar u^{\sigma,\alpha}(p),
\]
with the same pole residue after the phase convention is fixed.

The time-ordered two-point function then has the one-particle pole
\[
  \langle\Omega|
  \operatorname T\{
    \widehat\psi_\alpha(x)\widehat{\bar\psi}^{\beta}(y)
  \}
  |\Omega\rangle
  =
  Z_\psi
  (-\not\partial_x+m)_\alpha{}^\beta
  \Delta_F(x-y;m^2)
  +
  G_{\alpha}{}^{\beta,\mathrm{cont}}(x-y),
\]
where \(G^{\mathrm{cont}}\) is the contribution from states outside the
isolated one-particle mass shell. In momentum space the pole term is
\[
  Z_\psi
  \frac{-\ii(-\ii\not k+m)_\alpha{}^\beta}
       {k^2+m^2-\ii\epsilon}.
\]

Spinorial LSZ reduction applies the same Haag--Ruelle and pole-extraction
logic as the scalar case. The difference is that the residue has spinor
indices. Amputation of an external spinor leg therefore consists of:
\[
  \text{multiply by }(k^2+m^2),\qquad
  \text{take }k^0\to\omega_{\vec k},\qquad
  \text{contract with }u^\sigma(p)\text{ or }v^\sigma(p),
\]
with the choice of \(u\) or \(v\) determined by the external charge and by
whether the particle is incoming or outgoing. Thus external spinor factors in
perturbative scattering are the one-particle pole intertwiners of local
spinor fields.

The next chapter turns to massless particles. The logical input will again be
the one-particle representation theory, now with the massless little group and
helicity replacing the massive \(SU(2)\) spin label.
